\subsection{Damping-Like and Field-Like Spin--Orbit Torque}

\paragraph{Equation.}
With a spin-polarization direction $\unitvect{\sigma}$, the spin--orbit torque can be
written as
\begin{equation}
    \tauSOT
    = \tau_{\mathrm{DL}}\,
      \mvec\times(\unitvect{\sigma}\times\mvec)
      + \tau_{\mathrm{FL}}\,
      \mvec\times\unitvect{\sigma}.
    \label{eq:sot-dl-fl}
\end{equation}
It enters the Gilbert equation additively:
\begin{equation}
    \frac{\dd\mvec}{\dd t}
    = -\gamma_0\,\mvec\times\Heff
      + \alpha\,\mvec\times\frac{\dd\mvec}{\dd t}
      + \tauSOT.
    \label{eq:llg-sot}
\end{equation}

\paragraph{Definitions.}
$\tau_{\mathrm{DL}}$ and $\tau_{\mathrm{FL}}$ are torque amplitudes with units of
$\mathrm{s^{-1}}$ in this normalization.  The unit vector $\unitvect{\sigma}$ gives
the injected spin-polarization direction.

\paragraph{Assumptions.}
The torque is spatially uniform and the magnetization remains normalized.

\paragraph{Sign convention.}
The identities
$\mvec\times(\unitvect{\sigma}\times\mvec)
=-\mvec\times(\mvec\times\unitvect{\sigma})$ are equivalent, but their scalar
coefficients have opposite signs.  Record the current direction, electron-flow
direction, layer normal, and definition of $\unitvect{\sigma}$ when assigning torque
amplitudes.

\paragraph{Source.}
Add the source for the torque convention and amplitude model here.

\paragraph{Implementation notes.}
If the numerical solver converts the Gilbert equation to an explicit form, it must
also transform the torque terms consistently; do not only divide Eq.~\eqref{eq:llg-sot}
by $1+\alpha^2$.
